\subsection{Проектирование системы}

\subsubsection{Структура системы}

Разрабатываемая программная система имеет клиент-серверную структуру, соответствующую трёхзвенной модели: уровень представления (клиентские веб- и мобильные приложения), уровень приложений (серверная часть с программным интерфейсом и бизнес-логикой), уровень данных (реляционная СУБД PostgreSQL).

В состав системы входят следующие компоненты:
\begin{enumerate}
    \item мобильное приложение для жителей~--- кроссплатформенное приложение на React Native, обеспечивающее фиксацию проблем городской инфраструктуры, просмотр обращений на интерактивной карте и отслеживание статусов;
    \item диспетчерская веб-панель~--- автоматизированное рабочее место диспетчера, реализованное на React с библиотекой Яндекс.Карты;
    \item серверная часть~--- программный интерфейс на платформе ASP.NET Core, обеспечивающий обработку бизнес-логики и хранение данных в СУБД PostgreSQL с расширением PostGIS.
\end{enumerate}

Клиентские приложения взаимодействуют с серверной частью посредством протокола HTTPS, обмениваясь данными в формате JSON. На рисунке~\ref{fig:architecture} представлена диаграмма компонентов программной системы.

\begin{figure}[H]
    \centering
    \includegraphics[width=\textwidth]{architecture.png}
    \caption{Диаграмма компонентов программной системы}
    \label{fig:architecture}
\end{figure}

\subsubsection{Структура серверной части}

Серверная часть построена на основе чистой структуры в сочетании с подходом модульного монолита. На рисунке~\ref{fig:server-components} представлена структура серверного приложения, построенная на принципах чистой структуры с разделением на слои API, Application, Domain и Persistence.

\begin{figure}[H]
    \centering
    \includegraphics[width=\textwidth]{server-components.png}
    \caption{Структура серверного приложения}
    \label{fig:server-components}
\end{figure}

Серверное приложение включает следующие предметные модули:
\begin{enumerate}
    \item модуль пользователей~--- пользователи, роли, аутентификация, профиль;
    \item модуль обращений~--- приём и обработка обращений жителей, фотофиксация, привязка к геолокации, дедупликация в радиусе 50~м и отслеживание статусов;
    \item модуль событий~--- городские события с категоризацией, тегами и модерацией;
    \item модуль уведомлений~--- уведомления о смене статусов обращений;
    \item модуль файлов~--- загрузка и хранение медиафайлов.
\end{enumerate}

Структура каждого модуля организована по четырём слоям: предметный слой (бизнес-сущности, объекты-значения), прикладной слой (интерфейсы репозиториев и сервисы бизнес-логики), слой доступа к данным (реализация репозиториев через Entity Framework Core), слой программного интерфейса (контроллеры REST API).

\subsubsection{Проектирование базы данных}

Используется реляционная СУБД PostgreSQL с расширением PostGIS для геопространственных запросов. На рисунке~\ref{fig:erd} представлена ER-диаграмма базы данных программной системы.

\begin{figure}[H]
    \centering
    \includegraphics[width=0.85\textwidth]{erd-smartcity.png}
    \caption{ER-диаграмма базы данных программной системы}
    \label{fig:erd}
\end{figure}

Ключевые сущности модуля обращений:
\begin{itemize}
    \item Report~--- обращение жителя с описанием, координатами, статусом и ссылками на категорию и автора;
    \item ReportCategory~--- классификация проблем (дороги, освещение, благоустройство, ЖКХ, экология);
    \item MunicipalService~--- справочник ответственных служб с контактными данными;
    \item ReportCategoryService~--- связующая таблица категорий и служб для маршрутизации обращений;
    \item ReportStatusHistory~--- хронологическая запись изменений статуса с указанием диспетчера и комментария;
    \item ReportMedia~--- метаданные и URL-адреса фотографий обращений;
    \item ReportConfirmation~--- подтверждение обращения другим жителем (<<Тоже вижу проблему>>).
\end{itemize}

Для геолокационных запросов координаты хранятся в формате, поддерживаемом расширением PostGIS. Поиск осуществляется функцией ST\_DWithin.

На рисунке~\ref{fig:state-appeal} представлена диаграмма состояний, описывающая жизненный цикл обращения о проблеме городской инфраструктуры.

\begin{figure}[H]
    \centering
    \includegraphics[width=0.45\textwidth]{state-appeal.png}
    \caption{Диаграмма состояний: жизненный цикл обращения}
    \label{fig:state-appeal}
\end{figure}

\subsubsection{Проектирование взаимодействия компонентов}

Для детализации взаимодействия между компонентами системы рассмотрены ключевые сценарии. На рисунке~\ref{fig:seq-create} представлена диаграмма последовательности для процесса создания обращения из мобильного приложения с учётом механизма дедупликации.

\begin{figure}[H]
    \centering
    \includegraphics[width=\textwidth]{seq-create-appeal-mobile.png}
    \caption{Диаграмма последовательности: создание обращения из мобильного приложения}
    \label{fig:seq-create}
\end{figure}

На рисунке~\ref{fig:seq-dispatcher} представлена диаграмма последовательности для процесса обработки обращения диспетчером через веб-панель.

\begin{figure}[H]
    \centering
    \includegraphics[width=\textwidth]{seq-dispatcher.png}
    \caption{Диаграмма последовательности: обработка обращения диспетчером}
    \label{fig:seq-dispatcher}
\end{figure}

\subsubsection{Структура мобильного приложения}

Мобильное приложение разработано на платформе React Native с использованием TypeScript и инструментария Expo. На рисунке~\ref{fig:mobile-arch} представлена диаграмма структуры мобильного приложения.

\begin{figure}[H]
    \centering
    \includegraphics[width=\textwidth]{mobile-architecture.png}
    \caption{Структура мобильного приложения}
    \label{fig:mobile-arch}
\end{figure}

Структура мобильного приложения включает:
\begin{enumerate}
    \item экраны~--- визуальные компоненты пользовательского интерфейса: авторизация, карта, создание обращения, список обращений, карточка обращения, уведомления, профиль;
    \item навигация~--- модуль на основе expo-router с нижней панелью вкладок и стековым навигатором для модальных экранов;
    \item переиспользуемые компоненты~--- маркеры карты, карточки обращений, формы ввода, модальные окна, индикаторы загрузки;
    \item управление состоянием~--- React Context и пользовательские хуки, защищённое хранилище expo-secure-store для данных аутентификации;
    \item клиент API~--- модуль на основе Axios с аутентификацией на основе JWT, автоматическим обновлением токена и обработкой ошибок;
    \item сервисы устройства~--- модули доступа к камере и геолокации.
\end{enumerate}

На рисунке~\ref{fig:navigation} представлена диаграмма навигационной структуры мобильного приложения, отображающая связи между экранами и способы перехода.

\begin{figure}[H]
    \centering
    \includegraphics[width=\textwidth]{navigation-structure.png}
    \caption{Структура навигации мобильного приложения}
    \label{fig:navigation}
\end{figure}

\subsubsection{Проектирование пользовательского интерфейса мобильного приложения}

На основании требований к составу экранов были разработаны макеты ключевых экранов мобильного приложения. Экран авторизации (рисунок~\ref{fig:mockup-auth}) отображается при запуске для неаутентифицированного пользователя и содержит логотип, поля ввода и кнопку входа. Экран карты (рисунок~\ref{fig:mockup-map}) является центральным экраном и отображает городские события и обращения на интерактивной карте. Экран создания обращения (рисунок~\ref{fig:mockup-create}) предназначен для фиксации проблемы и включает выбор категории, описание, фотографию и указание местоположения.

\begin{figure}[H]
    \centering
    \includegraphics[width=0.3\textwidth]{screen-auth.png}
    \caption{Макет экрана авторизации}
    \label{fig:mockup-auth}
\end{figure}

\begin{figure}[H]
    \centering
    \includegraphics[width=0.3\textwidth]{screen-map.png}
    \caption{Макет экрана карты с событиями}
    \label{fig:mockup-map}
\end{figure}

\begin{figure}[H]
    \centering
    \includegraphics[width=0.3\textwidth]{screen-create-report.png}
    \caption{Макет экрана создания обращения}
    \label{fig:mockup-create}
\end{figure}

Экран списка обращений (рисунок~\ref{fig:mockup-list}) предназначен для просмотра обращений пользователя и включает вкладки фильтрации по статусу. Экран карточки обращения (рисунок~\ref{fig:mockup-detail}) отображает подробную информацию: фотографию, категорию, адрес, описание, счётчик подтверждений и хронологию статусов.

\begin{figure}[H]
    \centering
    \includegraphics[width=0.3\textwidth]{screen-reports-list.png}
    \caption{Макет экрана списка обращений}
    \label{fig:mockup-list}
\end{figure}

\begin{figure}[H]
    \centering
    \includegraphics[width=0.3\textwidth]{screen-report-detail.png}
    \caption{Макет экрана карточки обращения}
    \label{fig:mockup-detail}
\end{figure}

\begin{figure}[H]
    \centering
    \includegraphics[width=0.3\textwidth]{screen-notifications.png}
    \caption{Макет экрана уведомлений}
    \label{fig:mockup-notifications}
\end{figure}

\begin{figure}[H]
    \centering
    \includegraphics[width=0.3\textwidth]{screen-profile.png}
    \caption{Макет экрана профиля пользователя}
    \label{fig:mockup-profile}
\end{figure}
